\section{How Everything Merges Together: the Rational Hash and its Encoding in SmallWood}
\label{sec:Everything-Merges}

In this section, we will first describe the BBS‑style rational hash and its proof-friendly identity, then explain how witness rows are packed and committed, how each constraint family is built and wired, how batches $Q_i$ are formed and checked in the PIOP, and finally how the small‑field variant changes costs and soundness.

\subsection{Rational hash, target, and cleared identity}
\label{subsec:merge-rational}

\paragraph{Definition (hash and ISIS target).}
Public parameters fix $A\in R_q^{n\times m}$ and $(B_0,b_1)\in R_q^{n}\times R_q^{n}$. For a message $u$ and randomness $\chi=(x_0,x_1)$ we set
\[
  t \ :=\ h_{u,\chi}(B)\ =\ B_0\cdot(1;u;x_0)\ \Hdiv\ \bigl(b_1-\mathbf 1_n\!\cdot x_1\bigr)\ \in R_q^{n},
\]
and the preimage relation is $A s \equiv t \pmod q$ with $\|s\| \le \beta$

\paragraph{Guard (denominator invertibility).}
We enforce slot‑wise invertibility of $\bigl(b_1-\mathbf 1_n\!\cdot x_1\bigr)$ by resampling $x_1$ until every slot is invertible in $R_q$; this realizes the message‑subspace guard and legitimizes the cleared form below.

\paragraph{Exact evaluation (implementation contract).}
We compute $t$ in the evaluation/NTT domain and keep it there until needed downstream:
\begin{algorithm}[H]
\caption{\textsf{ComputeBBSHash}$(B_0,b_1;u,x_0,x_1)\rightarrow t$ (evaluation/NTT domain)}
\label{alg:compute-bbs-exact}
\begin{algorithmic}[1]
\State Lift $u,x_0,x_1$ to the evaluation domain.
\State $r \gets B_0^{(\mathrm{const})}\ +\ B_0^{(u)}\Had u\ +\ B_0^{(x)}\Had x_0$;\quad $d \gets b_1 - x_1$.
\State If any slot of $d$ is non‑invertible, resample $x_1$; else set $t \gets r \Had d^{-1}$.
\State \Return $t$.
\end{algorithmic}
\end{algorithm}

\paragraph{Cleared identity (degree‑$\le 2$).}
Writing $X_1=\mathbf 1_n\!\cdot x_1$, the equality $A s \equiv t$ is equivalent (under the guard) to
\begin{equation}\label{eq:cleared-bbs-merged}
  (b_1\Had A)\,s \;-\; (A s)\,X_1 \;-\; B_0\cdot(1;u;x_0)\ \equiv\ 0 \pmod q,
\end{equation}
which has total degree $\le 2$ in $(s,u,x_0,x_1)$. This is the main equality constraint we will encode in PACS/PIOP.

\subsection{Packing model and committed rows}
\label{subsec:merge-packing}


\paragraph{Grid and degrees.}
Let $\Omega=\{\omega_1,\dots,\omega_s\}\subset\F_q$ be the evaluation grid (take distinct NTT nodes, with $S_0:=\sum_{\omega\in\Omega}1=|\Omega|\not\equiv 0\pmod q$). Each committed row is interpolated to a polynomial with degree $\le s+\ell-1$ by fixing its values on $\Omega$ and $\ell$ blinded values on $S\subset\F_q\setminus\Omega$. DECS enforces the degree bound.

\paragraph{Packing policy.}
Signature rows $s$ use \emph{coefficient packing}: column $j$ carries the $j$‑th coefficient of the ring element. The message/randomness rows $(u,x_0)$ are \emph{column‑constant} (degree‑0) rows over $\Omega$. The scalar $x_1$ is carried by a column‑constant row $X_1$.

\paragraph{Witness blocks (to localize products).}
We form the per‑column blocks
\[
  w_1 := (s,u,x_0), \qquad w_2 := x_1, \qquad w_3 := w_1\cdot w_2,
\]
with the product taken slot‑wise. Denote by $W_1(X),W_2(X),W_3(X)$ the corresponding row‑polynomial stacks.

\paragraph{Public tables as row polynomials.}
The public linear maps $A$, $(b_1\!\Had A)$ and the routed coefficients from $B_0$ are similarly lifted to \emph{public} row polynomials: $A_j(X)$, $(b_1\!\Had A)_j(X)$, and $B_{0,j}^{(\cdot)}(X)$.

\subsection{Constraint assortment and wiring}
\label{subsec:merge-constraints}

Now, let us describe the exact constraint families and show how they wire to \eqref{eq:cleared-bbs-merged} in order to have them verify the signature validity.

\paragraph{(C1) Parallel product gates (slot‑wise, quadratic).}
For every gated component $t$ in $w_1$,
\begin{equation}\label{eq:par-gate}
  F^{\mathrm{par}}_{t}(X)\ :=\ W^{(t)}_{3}(X)\;-\;W^{(t)}_{1}(X)\,W_{2}(X)\ \equiv\ 0
\end{equation}
pins $W_3=W_1\cdot W_2$ on $\Omega$ (and at fresh points via the PIOP tail).

\paragraph{(C2) Aggregated cleared identity (row‑polynomial form).}
For each public row index $j$,
\begin{equation}\label{eq:agg-bbs}
\begin{aligned}
F^{\mathrm{bbs}}_j(X)
  :=\;& \underbrace{\sum_{k\in\mathcal S} (b_1\!\Had A)_{j,k}(X)\,S_k(X)}_{\textstyle (b_1\!\Had A)s}  \\
  &-\underbrace{\sum_{k\in\mathcal S} A_{j,k}(X)\,W^{(k)}_{3}(X)}_{\textstyle (A s)\,x_1}  \\
  &-\underbrace{\Bigl(B_{0,j}^{(0)}(X)
     +\!\!\sum_{t\in\mathcal U}\!\! B_{0,j}^{(u,t)}(X)\,U_t(X)
     +\!\!\sum_{\tau\in\mathcal X}\!\! B_{0,j}^{(x_0,\tau)}(X)\,X_{0,\tau}(X)\Bigr)}_{\textstyle B_0\cdot(1;u;x_0)}
  \equiv 0 .
\end{aligned}
\end{equation}

Here $\mathcal S$ indexes signature rows, $\mathcal U$ message rows, and $\mathcal X$ the $x_0$ rows. This is the packed analogue of \eqref{eq:cleared-bbs-merged}.

\paragraph{(C3) Alphabet membership for $(u,x_0,x_1)$ (parallel, high‑degree but bounded).}
For each $P\in\{u,x_0,x_1\}$ we enforce a single vanishing‑polynomial row
\[
  P_B\!\big(P(\cdot)\big)\ \equiv\ 0 \quad\text{on }\Omega,\qquad
  P_B(X)=\prod_{i=-B}^{B}\bigl(X-\langle i\rangle_q\bigr).
\]

\paragraph{(C4) $\ell_\infty$ chain for signature rows (parallel).}
For each signature row and each column value $v$ we write a balanced‑radix decomposition with digit membership, carries, and recomposition:
\[
  P_{\mathcal D}(d_{k})=0,\qquad
  R\,d_{k}+c_{k}=d'_{k}+c_{k+1},\qquad
  v=\sum_{k=0}^{L-1}R^k d_k,
\]
with $R=2^W$ and balanced digit sets. The digit vanishing degree drives the parallel‑degree parameter $d$, while carry/recomposition are quadratic. This certifies the $\ell_\infty$ bound used by your global $\ell_2$ budget.

\paragraph{Low‑degree vs.\ degree‑constrained.}
(C1)–(C2) are genuinely low‑degree; (C3)–(C4) are high‑degree but degree‑bounded. The maximal degree among parallel constraints determines the PIOP degree parameter $d_Q$.

\subsection{Batching, masks, and PIOP checks}
\label{subsec:merge-batching}

\paragraph{What the PIOP sees.}
All rows live under the same LVCS/DECS commitment. The verifier checks linear combinations at verifier‑chosen coordinates; PCS wraps these into point openings.

\paragraph{Batch polynomials.}
We form $\rho$ batches
\[
  Q_i(X)\ =\ M_i(X)\ +\ \sum_t \Gamma'_{i,t}\,F^{\mathrm{par}}_{t}(X)\ +\ \sum_j \gamma'_{i,j}\,F^{\mathrm{bbs}}_{j}(X),
\]
with Fiat–Shamir weights $(\Gamma',\gamma')$. Each $M_i$ is picked so that $\Sigma_\Omega Q_i=0$.

\paragraph{Mask construction (large‑field case).}
Let $S_k:=\sum_{\omega\in\Omega}\omega^k$. Choose random $a_1,\dots,a_{d_Q}$ and set
\[
  a_0 \ =\ -\,S_0^{-1}\,\Bigl[
     \sum_{k=1}^{d_Q} a_k\,S_k
     \;+\; \sum_t \Gamma'_{i,t}\,\Sigma_\Omega F^{\mathrm{par}}_{t}
     \;+\; \sum_j \gamma'_{i,j}\,\Sigma_\Omega F^{\mathrm{bbs}}_{j}
  \Bigr],
\]
then $M_i(X)=\sum_{k=0}^{d_Q} a_k X^k$. The verifier checks $Q_i$ on $\Omega$ and at fresh tail points and enforces $\Sigma_\Omega Q_i=0$.

\paragraph{Soundness knobs (pointer).}
The error splits into DECS degree, masked‑sum, tail, and LVCS tail‑binding components; the degree term uses the \emph{true} degrees from (C3)–(C4). Parameters instantiations will be discussed later.

\subsection{Small‑field variant}
\label{subsec:smallfield}


Let $K/\F$ be an extension of degree $\theta$. We run the PIOP over $K$ while keeping LVCS/DECS over $\F$.

\paragraph{How the interface works.}
A single evaluation in $K$ expands to its $\theta$ $\F$‑coordinates for LVCS queries (row/query replication by $\theta$). Batches live in $K[X]$, masks $M_i$ are built in $K[X]$ so that $\sum_{\omega\in\Omega} Q_i(\omega)=0$ holds \emph{in $K$}, and DECS enforces the same row‑degree bound as in the large‑field case.

\paragraph{Gains and trade‑offs.}
Because $|K|=|\F|^\theta$, we can lower both the amount of tail points $\ell'$ tail point and masked‑sum batch $\rho$  while preserving the target soundness (PIOP error $\le (1+d_Q)/|K|$). This reduces challenge material. The trade‑off is a factor‑$\theta$ replication of LVCS queries/rows (over $\F$), but the total opening count still drops when $\theta$ is chosen to meet the security level.
